\documentclass[12pt]{article}

\usepackage[margin=1in]{geometry}
\usepackage{graphicx}
\usepackage{booktabs}
\usepackage{amsmath}
\usepackage{natbib}
\usepackage{setspace}
\usepackage{hyperref}

\onehalfspacing

\title{Replicating Table 3 (Fivefold LATE) from Chernozhukov et al.\ (2018)}
\author{Seth Abramowitz}
\date{\today}

\begin{document}

\maketitle

\begin{center}
\small GitHub Repository: \href{https://github.com/Seth-Ab/552-Replication}{https://github.com/Seth-Ab/552-Replication}
\end{center}

\section{The Original Paper}

\citet{chernozhukov2018} studies how Double/Debiased Machine Learning can be
used to estimate causal parameters while allowing flexible machine-learning
methods in the nuisance-function stage. The paper's central concern is whether
modern prediction tools can be combined with orthogonal score construction and
cross-fitting in a way that preserves valid inference for treatment effects.

The application I replicate uses the 1991 SIPP 401(k) data. In that exercise,
the paper estimates the effect of 401(k) participation on net financial assets
while treating participation as endogenous and using 401(k) eligibility as an
instrument. That setting makes Table 3 a natural illustration of the
interactive instrumental-variables version of Double Machine Learning.

\section{The Result of Interest}

I replicate the fivefold cross-fitting row of Table 3 from
\citet{chernozhukov2018}. More specifically, I target the row reporting the
Local Average Treatment Effect of 401(k) participation on net financial assets
across the paper's set of nuisance learners. This is the most useful result to
reproduce because it is the paper's main empirical IV example and because it
shows whether the qualitative conclusion is robust to different first-stage
machine-learning choices.

\section{Data and Methods}

The raw data are stored in \texttt{input/sipp1991.dta}, the 1991 SIPP household
sample used in the DoubleML 401(k) application. The variables used downstream
are documented in \texttt{input/data\_dictionary.md}. The key variables are net
financial assets as the outcome, 401(k) participation as the treatment, 401(k)
eligibility as the instrument, and a set of demographic and financial controls.

The replication pipeline is split into preprocessing and analysis. The script
\texttt{code/preprocess.py} reads the raw Stata file from \texttt{input/}, keeps
the variables used in the replication, drops missing values, and writes a
deterministic analysis-ready file to \texttt{temp/clean\_data.csv}. The script
\texttt{code/analysis.py} then reads that cleaned file, estimates the IIVM
specification with fivefold cross-fitting, writes the tracked output tables into
\texttt{output/tables/}, and prints a console comparison between the published
row and the replicated row.

\section{Replication Results}

\input{../output/tables/main_result.tex}

The generated table reports the published targets in Panel A and the replication
results in Panel B. The main qualitative pattern carries over: the estimated
LATE remains positive across learners, and the broad magnitudes are comparable
to the original table. The replication therefore succeeds at reproducing the
paper's central empirical message even though it does not match every learner
exactly.

The remaining discrepancies are plausible given the structure of the exercise.
Learner-specific tuning choices can move nuisance predictions substantially,
especially for penalized and neural-network specifications. In addition, even
when the estimand and cross-fitting structure are aligned, implementation
details in repeated sample splitting and aggregation can affect the final row.

\section{What I Learned}

The easiest part of the replication was identifying the target result once I
committed to reproducing only one row instead of the entire paper. The hardest
part was making the machine-learning nuisance stage stable and reproducible
across several learners while keeping the full pipeline automated from raw data
through the final PDF.

The main lesson from the exercise is that reproducibility depends as much on
workflow discipline as it does on the estimator itself. A clear directory
structure, deterministic preprocessing, generated LaTeX fragments, and a real
Makefile DAG make it much easier for another researcher to understand what was
done and to rerun the analysis without guessing. If I were the original author,
the most helpful additional detail would be more explicit documentation of the
learner tuning and repeated-split aggregation choices used for the final table.

\bibliographystyle{apalike}
\bibliography{references}

\end{document}
